\documentclass{article}
\usepackage{graphicx}
\usepackage{amsmath}
\usepackage{amssymb}
\usepackage{mathtools}
\usepackage{amsthm}
\usepackage{geometry}
\usepackage{xcolor}
\usepackage{booktabs}
\usepackage{hyperref}
\geometry{margin=0.75in}
\setlength{\parindent}{0pt}

% ---- macros (matching javan_theory/macros.tex where sensible) ----
\newcommand{\note}[1]{\textcolor{cyan}{\textbf{**#1}}}
\newcommand{\xb}{\boldsymbol{x}}
\newcommand{\pb}{\boldsymbol{p}}
\newcommand{\vb}{\boldsymbol{v}}
\newcommand{\qb}{\boldsymbol{q}}
\newcommand{\kb}{\boldsymbol{k}}
\newcommand{\eb}{\boldsymbol{e}}
\newcommand{\gb}{\boldsymbol{g}}
\newcommand{\ub}{\boldsymbol{u}}
\newcommand{\mb}{\boldsymbol{m}}
\newcommand{\rb}{\boldsymbol{r}}
\newcommand{\db}{\boldsymbol{d}}
\newcommand{\cb}{\boldsymbol{c}}
\newcommand{\Xb}{\boldsymbol{X}}
\newcommand{\Ab}{\boldsymbol{A}}
\newcommand{\Bb}{\boldsymbol{B}}
\newcommand{\Cb}{\boldsymbol{C}}
\newcommand{\Db}{\boldsymbol{D}}
\newcommand{\Gb}{\boldsymbol{G}}
\newcommand{\Ib}{\boldsymbol{I}}
\newcommand{\Jb}{\boldsymbol{J}}
\newcommand{\Kb}{\boldsymbol{K}}
\newcommand{\Mb}{\boldsymbol{M}}
\newcommand{\Pb}{\boldsymbol{P}}
\newcommand{\Rb}{\boldsymbol{R}}
\newcommand{\Sb}{\boldsymbol{S}}
\newcommand{\Tb}{\boldsymbol{T}}
\newcommand{\Vb}{\boldsymbol{V}}
\newcommand{\Wb}{\boldsymbol{W}}
\newcommand{\Yb}{\boldsymbol{Y}}
\newcommand{\etab}{\boldsymbol{\eta}}
\newcommand{\pib}{\boldsymbol{\pi}}
\newcommand{\omegab}{\boldsymbol{\omega}}
\newcommand{\one}{\boldsymbol{1}}
\newcommand{\normsq}[2]{\Vert {#1} \Vert_{#2}^2}
\newcommand{\tr}{\textrm{tr}}
\newcommand{\softmax}[1]{\textrm{softmax}(#1)}
\newcommand{\We}{\boldsymbol{W}_E}
\newcommand{\Wu}{\boldsymbol{W}_U}
\newcommand{\Wov}{\boldsymbol{W}_{OV}}
\newcommand{\lattn}{\mathcal{L}_{\text{attn}}}
\newcommand{\Tx}{\Tb^{(x)}}
\newcommand{\Tz}[1]{\Tb^{(#1)}}
\newcommand{\Tcond}[1]{\Tb^{|#1}}
\newcommand{\Btau}{\Bb^{(\tau)}}
\newcommand{\vocab}{{\vert \mathcal{V}\vert}}
\newcommand{\state}{{\vert \mathcal{S}\vert}}
\newcommand{\alphatil}{\Tilde{\alpha}}
\newcommand{\gammatil}{\Tilde{\gamma}}
\newcommand{\RF}{\mathrm{RF}}

\theoremstyle{definition}
\newtheorem{definition}{Definition}[section]
\newtheorem{prop}{Proposition}
\newtheorem{lemma}{Lemma}
\newtheorem{remark}{Remark}
\newtheorem{theorem}{Theorem}[section]
\newtheorem{corollary}{Corollary}

\title{Feature-resolved attention in one-layer transformers:\\
attribution and intervention in the constrained belief-update framework}
\author{Sprint worker B (run1\_twin)}
\date{July 15, 2026}

\begin{document}
\maketitle

\begin{abstract}
Feature-resolved attention (FRA) decomposes a head's attention scores (QK) and
transported content (OV) through a sparse feature basis. Empirically, FRA edits
control association-specific attention behaviors but fail on aggregated,
broad concepts. We derive this boundary inside the constrained
belief-update theory of one-layer transformers. Four results. (1) A
\emph{gauge theorem}: whenever the attention pattern is token-independent
across sequences, every content-involving component of the attention scores
that depends on token identity is either a softmax row constant or a
hard-cancelled channel pair;
feature-resolved score mass is then gauge, not mechanism, and full key-side
FRA-QK cuts leave the pattern exactly unchanged. Trained models satisfy the
premise only in expectation --- they carry a small, loss-bearing
content-gating (worth $2\times10^{-4}$ nats on Mess3, about $1\%$ of the
attention benefit, and not induction-shaped) --- so the theorem's role is to expose score mass as
causally empty by default, with partial-cut effects predicted to four
decimals by a one-parameter common-mode formula. (2) A closed form for FRA-OV attributions at the theoretical
weights: $\zeta^{d-s}\,\gb(z_s)$, identifying FRA-OV with effective subspace
attention in an SAE basis, exact up to a derived softmax normalization warp
$(1-\zeta^d)^{-1}$. (3) A \emph{quadratic removal law}: near a Bayes-optimal
model, the behavioral removal fraction of an FRA-OV cut at gain $c$ on a path
of share $\rho$ is $\RF(c)\approx(c\rho)^2$; honest severing ($c=1$) of a
quarter-share path removes $6\%$ of behavior, and the gain-tuned null sits at
$c^\ast = 1/\rho$. One fitted parameter reproduces the fra\_hmm\_toy gain
sweeps, including $c^\ast\approx 4$. (4) Use/presence as rank conditions:
a single gain can null the readout projection (use) exactly when all paths
aggregate the same statistic, but can null the full concept Jacobian
(presence) only at a rank bottleneck; path cuts therefore preserve
decodability while content cuts at the carrier are the only presence
reducers, with collateral bounded by an entanglement tax.
\end{abstract}

\section{Setup}

\subsection{Process and the constrained belief update}

Mess3 is an edge-emitting hidden Markov process with three hidden states and
three tokens, parameterized by an emission fidelity $\alpha$ and a transition
rate $x$; writing $\beta = (1-\alpha)/2$ and $y = 1-2x$, the token-labeled
transition matrices are
\begin{equation}
\Tz{0} =
\begin{pmatrix}
\alpha y & \beta x & \beta x \\
\alpha x & \beta y & \beta x \\
\alpha x & \beta x & \beta y
\end{pmatrix},
\end{equation}
with $\Tz{1}, \Tz{2}$ the corresponding matrices with the $\alpha$-column at
positions $2, 3$. The marginal transition matrix $\Tb = \sum_z \Tz{z}$ has
stationary distribution $\pib = (\tfrac13,\tfrac13,\tfrac13)$ and eigenvalues
$\{1, \zeta, \zeta\}$ with $\zeta = 1-3x$; the spectral projector onto the
decaying eigenspace is $\Pb_\zeta = \Ib - \one\pib$. Writing
$\Tcond{z}$ for the row-normalization of $\Tz{z}$, define the token
displacement
\begin{equation}\label{gdef}
\gb(z) = \pib\, \Tcond{z} - \pib,
\end{equation}
a row vector in the simplex plane with $\sum_z \gb(z) = 0$. The
\textit{constrained belief update} of Piotrowski et al.\ is the best
approximation to the Bayesian belief state at position $d$ expressible as a
sum of independent per-source corrections,
\begin{equation}\label{constrained}
\rb_d = \pib + \sum_{s=1}^{d} \zeta^{\,d-s}\, \gb(z_s)
=: \pib + \ub_d,
\end{equation}
where all lag dependence sits in the scalar $\zeta^{d-s}$ and all token
dependence in the displacement $\gb(z_s)$. For our reference parameters
$(\alpha, x) = (0.6, 0.1)$, $\zeta = 0.7$, the cross entropy ladder is
$1.0809$ (Bayes) $\le 1.0833$ (constrained) $\le 1.0986$ (context-free
stationary predictor), so the constrained update captures $86\%$ of the
available predictive information. (A convention note that applies
throughout: cross entropies are Monte Carlo estimates on a fixed evaluation
set, and different experiments use different sets and position ranges, with
between-set offsets of order $10^{-3}$; we compare CE values only within
one set, and quote each experiment's own references.)

\subsection{Transformer, feature basis, and the FRA objects}

A one-layer attention-only transformer processes token sequences with
residual input $\xb_s = \eb(z_s) + \pb_s$ (token embedding plus additive
learned positional embedding, no BOS token, no layer normalization in our
reference configuration). A causal softmax head has scores
$\sigma(d,s) = \qb_d \cdot \kb_s/\sqrt{d_h}$ bilinear in the residuals,
pattern $\Ab(d,s) = \softmax{\sigma(d,\cdot)}$ over $s \le d$, and value
vectors $\vb_s = \Wov\, \xb_s$; a linear readout $f$ maps the attention
output $\sum_{s\le d}\Ab(d,s)\vb_s$ to the simplex plane.

A sparse dictionary (for us, a TopK SAE with decoder directions $\db_j$ and
activations $z_j(s)$) decomposes the residual at each position exactly into
active features plus a decoder bias and an error term. Substituting this
decomposition into the bilinear score and into the value path makes both
exactly feature-resolved:
\begin{align}
\sigma(d,s) &= \sum_{a, b} \sigma_{ab}(d,s), &
\sigma_{ab}(d,s) &= \frac{(\phi(t_a(d))\Wb_Q)\cdot(\phi(t_b(s))\Wb_K)}{\sqrt{d_h}},
\label{fraqk}\\
\cb_d &= \sum_{s \le d} \Ab(d,s) \sum_b \phi(t_b(s))\Wov,
\label{fraov}
\end{align}
where $t_a, t_b$ range over the per-side terms (active features, decoder
bias, error, constant) and $\phi$ is the identity (or the frozen
per-position layer-norm map when the model has LN). \textit{FRA-QK
attributions} are the score terms $\sigma_{ab}$ with a feature on the query
or key side; \textit{FRA-OV attributions} are the per-feature transported
contributions in \eqref{fraov}. The corresponding \textit{cuts} subtract $c$
times the attribution of a chosen feature set $S$ from the scores or from
the attention output; $c = 1$ is exact severing of the decomposition term,
$c > 1$ injects a counter-signal of the same shape.

\section{Realization: exact softmax construction and the normalization warp}

The first question is whether the constrained update \eqref{constrained} is
realizable at all by a single softmax head, since softmax rows sum to one
while the target kernel mass $\sum_{s\le d}\zeta^{d-s} = (1-\zeta^d)/(1-\zeta)$
grows with $d$.

\begin{prop}[exact realization]\label{prop:realization}
Let the scores be key-side only, $\sigma(d,s) = s\,(-\ln\zeta) + b(s)$ with
$b(1) = 0$ and $b(s) = \ln(1-\zeta)$ for $s \ge 2$. Then the causal softmax
pattern is
\begin{equation}\label{exactpattern}
\Ab(d,1) = \zeta^{\,d-1}, \qquad
\Ab(d,s) = (1-\zeta)\,\zeta^{\,d-s} \quad (2 \le s \le d),
\end{equation}
and with values $f(\vb_1) = \gb(z_1)$, $f(\vb_s) = \gb(z_s)/(1-\zeta)$ for
$s \ge 2$, the attention output equals $\ub_d$ exactly, at every position
and for every sequence.
\end{prop}

Row-stochasticity is a one-line telescoping check:
$\zeta^{d-1} + (1-\zeta)\sum_{k=0}^{d-2}\zeta^k = 1$. Position $1$ acts as
the attention sink in the absence of a BOS token: it absorbs the excess
softmax mass $\zeta^{d-1}$ (inflated by $1/(1-\zeta)$ relative to the pure
geometric profile), compensated by a value vector smaller by exactly
$(1-\zeta)$. This reproduces, as an exact statement, the empirical
observations of Piotrowski et al.\ that $A_{2,1}$ is anomalously large and
$\vb_1$ anomalously small.

\begin{prop}[normalization warp]\label{prop:warp}
Suppose instead the value scale is position-independent,
$f(\vb_s) = c\,\gb(z_s)$ --- which is what additive positional embeddings
force, since $\Wov\,\eb(z)$ cannot vary with $s$. Then no scores, even
token-dependent ones, realize $\ub_d$ exactly for all $d$: on a
constant-token sequence the output is $c\,\gb(z)$ while the target is
$\frac{1-\zeta^d}{1-\zeta}\gb(z)$, which forces $c$ to depend on $d$. For
the natural linear scores $\sigma = s(-\ln\zeta)$ the output equals
$\lambda_d\, \ub_d$ with
\begin{equation}\label{warp}
\lambda_d = \frac{c\,(1-\zeta)}{1-\zeta^{\,d}},
\end{equation}
a token-independent multiplicative warp. The choice $c = 1/(1-\zeta)$ is the
unique one that is asymptotically exact, leaving relative overshoot
$\zeta^d/(1-\zeta^d)$, geometric in $d$ and largest at early positions.
Token-independent positional value components cannot repair the warp: they
add the same vector to every sequence of length $d$, while the error
$(\lambda_d - 1)\ub_d$ is token-dependent; indeed since $E[\ub_d] = 0$ the
mean-square-optimal positional correction is exactly zero.
\end{prop}

Both propositions were verified to float precision ($2.8\times10^{-16}$) by
direct construction. A trained one-layer model (Adam, cross entropy, $96.5\%$
of the Bayes gap closed, evaluated on positions $\ge 2$) \emph{carries} the warp: the per-position scalar
$w(d)$ relating its readout to $\ub_d$ follows $1/(1-\zeta^d)$ with damped
amplitude ($w(1) = 2.67$ against the predicted $3.33$, i.e.\ $\approx 80\%$;
RMSE $0.15$ against the warp profile
versus $0.32$ against a flat profile). We propose this as the quantitative
content of the previously unexplained ``scalar discrepancy of the first two
embedding vectors'' reported in the constrained-belief paper.

\subsection{The profile cross entropy actually prefers}

The ansatz fixes the lag kernel to $\zeta^k$ --- the belief-decay rate of
the process. Optimizing the kernel itself inside the predictor family
$\rb_d = \pib + \sum_k w(k)\, \gb(z_{d-k})$ (Adam on Monte Carlo cross
entropy; evaluation on a common held-out set) shows the objective wants
something systematically different: the optimal kernel boosts the diagonal
($w(0) = 1.69$ against the ansatz's $1$) and decays \emph{faster} than the
process ($\text{rate} \approx 0.61 < \zeta = 0.7$), reaching CE $1.08018$
against $1.08071$ for the best gain-scaled $\zeta^k$ and $1.07959$ for
Bayes. The retrained positional-only softmax transformer lands on the same
choice (fitted late-row rate $0.575$). The mechanism is the
partial-regression shrinkage of the one-layer MSE theory: additive
per-source evidence over an autocorrelated sequence double-counts, so
optimal weights shrink the correlated tail relative to $\zeta^k$ and keep
the exact lag-$0$ evidence --- the cross-entropy analogue, in direction and
roughly in size, of that theory's prediction that optimal profile rates sit
below the process eigenvalue. Two further checks: releasing the kernel to
vary freely per destination row gains nothing ($1.08018$, a numerical
cross-entropy analogue of Toeplitz optimality), and the free row kernels
show \emph{no} early-row boost (relative lag-1 gain $\approx 1.0$ at all
rows against the warp's $1.96$ at $d=1$) --- the warp of
Proposition~\ref{prop:warp} is forced by softmax normalization, not
preferred by the objective.

\section{The QK gauge theorem}

The trained pattern of the one-layer model is token-independent (relative
standard deviation across token identities at fixed positions: $4$--$8\%$),
as the constrained-update theory requires. The following theorem
characterizes \emph{every} bilinear score function consistent with that
property; it is the reason feature-resolved score attributions carry no
mechanism.

\begin{theorem}[QK gauge]\label{thm:gauge}
Let $\sigma(d,s) = (\eb(z_d)+\pb_d)^\top \Mb\, (\eb(z_s)+\pb_s)$ and suppose
the causal-softmax pattern is the same for all token sequences (all token
values occurring with positive probability at every position). Decompose
$\sigma$ into content and position channels on each side. Then
\begin{enumerate}
\item content$\times$content is additively separable:
$\eb(w)^\top\Mb\,\eb(z) = a(w) + b(z)$;
\item content$\times$position is separable:
$\eb(w)^\top\Mb\,\pb_s = e(w) + f(s)$;
\item position$\times$content is separable,
$\pb_d^\top\Mb\,\eb(z) = u(d) + v(z)$, and the two key-side content parts
cancel: $b(z) + v(z) = \kappa$, a constant in $z$ --- a hard constraint, not
a softmax gauge;
\item consequently
$\sigma(d,s) = g(d, z_d) + f(s) + \pb_d^\top\Mb\,\pb_s$, where
$g = a + e + u + \kappa$ is a row constant. The pattern is
$\softmax{f(s) + \pb_d^\top\Mb\,\pb_s}$: every content-involving component
of the score is either a softmax row constant or part of a cancelling pair,
and the mechanism is carried by $f(s) + \pb_d^\top\Mb\,\pb_s$ alone.
\end{enumerate}
\end{theorem}

The proof varies one source token at a time (a single-entry change of a
softmax row must vanish), then the destination token (a row change must be
constant across the row). Three corollaries govern FRA-QK edits.

\begin{corollary}[full key-side cuts are exactly null]\label{cor:fullcut}
A key-side FRA-QK cut whose feature set covers the content at every position
subtracts $q_d^\top\Mb\,\eb(z_s) = \varphi(d, z_d)$ --- by the theorem the
same value at every $s$ --- i.e.\ a row constant. The pattern is exactly
unchanged at every gain.
\end{corollary}

\begin{corollary}[partial cuts measure gauge]\label{cor:partial}
Cutting a strict subset $S$ of tokens on the key side changes the pattern to
\begin{equation}\label{partialcut}
\Ab'(d,s) = \frac{\Ab(d,s)\, e^{-c\varphi\, 1_S(s)}}{1 - m_d + m_d\,
e^{-c\varphi}}, \qquad m_d = \textstyle\sum_{s \le d,\, z_s \in S} \Ab(d,s),
\end{equation}
with first-order row total variation $|c\varphi|\, m_d(1-m_d)$. The
controlling scalar $\varphi$ enters the clean model only as a row constant,
invisible to the loss: it is a gauge degree of freedom set by initialization
and training dynamics, so the size of a partial FRA-QK effect measures
gauge, not mechanism.
\end{corollary}

\begin{corollary}[query/key asymmetry]\label{cor:asym}
The mirrored full query-side cut subtracts $c[a(z_d) + e(z_d)]$ (gauge)
\emph{plus} $c[b(z_s) + f(s)]$: it re-exposes the cancelled key-content
variation $b(z_s)$ and removes the share of the positional mechanism $f(s)$
that the model implements through the content$\times$position channel.
Query-side cuts are therefore not null even at optimal weights.
\end{corollary}

The asymmetry matters in practice: in our trained one-layer model the lag
profile is carried dominantly by the content$\times$position channel through
the embedding \emph{common mode} (key-position main effect RMS $2.12$,
versus $1.04$ for the row-demeaned position$\times$position channel and
$0.014$ for the query-token main effect), so
Corollary~\ref{cor:asym} predicts visible query-side cut effects alongside
near-null key-side full cuts.

\textbf{Empirical scope of the theorem.} Cutting experiments on the trained
model draw the boundary precisely. Corollary~\ref{cor:partial} verifies
essentially exactly in an oracle content basis: the common-mode scalar is
row-constant to three decimals ($\varphi = 0.323$), and formula
\eqref{partialcut} predicts the measured pattern total variation and even
the loss change to four--five decimals across gains.
Corollary~\ref{cor:asym} verifies with the predicted sign and size
(query-side full cut costs $2.2\times$ the key-side cut at $c=1$). But
Corollary~\ref{cor:fullcut} \emph{fails} on the trained model, for an
instructive reason: the theorem's premise holds only in expectation. The
query-averaged pattern is token-independent (as originally reported), yet
flipping a single source token moves that source's attention weight by
$55\%$ relative on average (median $52\%$, seed 42) --- the trained model
genuinely content-gates. Dissecting
the gating pins down its worth and shape. Retraining the model with the
content$\to$QK path severed (positional scores kept free) yields a
cross-entropy ladder, per token on evaluation rows $1$--$30$:
\[
\underbrace{1.07898}_{\text{Bayes}} <
\underbrace{1.07966}_{\text{trained}} <
\underbrace{1.07981}_{\substack{\text{positional-only}\\ +\,\text{optimal }3\times3\text{ gate}}} <
\underbrace{1.07986}_{\text{positional-only}} <
\underbrace{1.08023}_{\substack{\text{best gain-scaled}\\ \text{constrained ansatz}}} <
\underbrace{1.08195}_{\text{constrained }\gamma=1}.
\]
Three conclusions. First, the model's advantage over the gain-optimized
ansatz is mostly a \emph{better positional profile} ($64\%$ of the
$5.7\times10^{-4}$-nat edge); the retrained value of content in QK is only
$2.0\times10^{-4}$ nats. Second, \emph{ablation is not value}: severing the
key-side content channel in place costs $0.0034$ nats, $17\times$ the
retrained value, because the model builds row-constant gauge structure on
the content channel and the in-place cut breaks that too. Third, the
learned content$\times$content interaction has an almost pure
\emph{same-token} shape ($0.9985$ of its Frobenius mass; seed-stable at
cosine $0.996$--$0.998$ across four seeds), yet a same-token
(induction-direction) gate is exactly null at the trained optimum (a tuned
gate gains $6\times10^{-7}$ nats, and the local gradient equals that of a
content-blind control), and the CE-optimal token-pair gate on the frozen
positional model has a \emph{different} shape (dominantly cyclic, cosine
$0.52$ with the learned interaction) and recovers only a quarter of the
content value. SGD thus reliably learns an induction-\emph{looking} score
structure whose causal content is nearly nil --- a decoy for any
score-mass reading of FRA-QK. The resulting trichotomy of FRA-QK score
mass: (a) row-constant gauge --- dominant and causally empty; (b) genuine
content-gating --- small where the data barely rewards it (here, worth
$\sim 1\%$ of the attention benefit), large where the data demands it
(induction regimes); (c) in the aggregation regime, pattern perturbations
that are behaviorally null by robustness of counting (next sections).
Attribution screens based on score mass see (a) and are fooled; only causal
tests isolate (b).

We closed the loop on (b) with the minimal induction toy: random tokens
carrying a per-sequence copy rule ($A$ is followed by $B$) planted at
random positions, plus a
previous-token input channel so a single head suffices. The trained score
interaction has the \emph{same} same-token shape as on Mess3 (Frobenius
share $0.999$ in both) --- and there it is the mechanism: the in-place
key-content cut at $c=1$ removes essentially the entire attention benefit
on rule positions ($\RF = 0.988$, $\Delta$CE $1.54$ nats against Mess3's
$0.0034$), severing content$\to$QK and retraining costs $0.90$ nats
($4400\times$ the Mess3 value of the identically-shaped structure), and a
single flipped source token collapses its attention weight by $98\%$.
The shape of a score matrix carries no information about its causal role;
the data distribution decides, exactly where the gauge theorem's premise
draws the line (\texttt{code/phaseD}). A methodological remark the
experiment produced on its own: with rule pairs planted at \emph{fixed}
positions, a positional shortcut absorbs most of the rule into the lag
profile and the measured content value drops threefold --- randomized
placement is part of the experimental design, not a detail.

\textbf{Basis quality gates the cuts --- and the requirement is
constructive.} A TopK SAE trained on the residual
stream splits each token into early/late \emph{position-band} latents
(token$\times$position mixtures). In that basis a ``content'' QK cut also
severs the positional lag mechanism the model runs on: removal fraction
$0.87$ at $c = 1$, versus $0.19$ for the oracle content cut. In an
entangled basis, FRA-QK edits measure basis entanglement --- neither gauge
nor concept. Repair experiments identify the property that fixes this. An
SAE trained on the position-demeaned residual restores the oracle number
exactly ($\RF = 0.185$ at reconstruction FVU $9\times10^{-5}$); an SAE
given position as a free side-input comes close ($0.216$); an $8\times$
larger dictionary does not ($0.344$) --- capacity alone leaves the
entanglement in place. The predictive statistic is the token-$R^2$ of the
\emph{summed} cut contribution $f_S$ ($1.000 / 0.972 / 0.872 / 0.830$
across the four bases, ordering the RF outcomes), while per-latent decoder
cosine to the token embeddings does not predict fidelity: the repaired
SAEs still score only $\approx 0.55$ per latent, because TopK distributes
each token's direction over several latents whose \emph{sum} is correct.
The faithfulness of an FRA cut is a property of the cut set, not of its
members.

Independently of any optimality assumption, cuts of \emph{slowly varying}
features obey a softmax stability bound: if the key-side score delta along a
row is $\delta(s) = \bar\delta + \varepsilon(s)$, the constant part cancels
in the softmax normalization at any magnitude, and the row total variation
of the pattern obeys $\mathrm{TV} \le \tfrac12\bigl(e^{\mathrm{osc}(\varepsilon)}
- 1\bigr)$ with $\mathrm{osc} = \max_s \varepsilon - \min_s \varepsilon$.
Per-sequence concept features (the mixture posterior below) have small
within-sequence variation, making their score contributions approximately
row-constant. On the trained mixture toy this bound is loose rather than
operative: the row-constant component carries $0.72$--$0.83$ of squared
score-delta mass under FRA-QK cuts, but the residual oscillations are
nat-scale and the pattern does move (median row TV $0.156$ at $c=1$); what
keeps the cuts behaviorally null there is aggregation robustness --- the
attention-weighted covariance between the score perturbation and the counted
tags is $\approx 0.02$, three orders below the perturbation scale.

\section{FRA-OV attributions and the intervention calculus}

\subsection{The attribution closed form}

At the weights of Proposition~\ref{prop:realization} (or \ref{prop:warp}),
the FRA-OV attribution of a token feature at source $s$ to the belief update
at destination $d$ is
\begin{equation}\label{ovclosed}
\Delta_{s\to d} = \Ab(d,s)\, f(\vb_s)
= \zeta^{\,d-s}\, \gb(z_s) \quad
\Bigl(\text{warped variant: } \tfrac{\zeta^{d-s}}{1-\zeta^d}\, \gb(z_s)\Bigr).
\end{equation}
This identifies FRA-OV, evaluated in an SAE basis whose token latents align
with the embedding directions, with the \emph{effective subspace attention}
of the attention-head-unit paper: the lag-resolved scalar $\zeta^{d-s}$
recovered by projecting transported content onto the known displacement
directions. The identification inherits that paper's gauge caveat: with
$H > 1$ heads only the head-sum is constrained. We tested this on the
negative-eigenvalue two-head mechanism ($\zeta = -0.5$, five seeds): the
head-sum effective attention collapses across seeds onto the oscillation
$\approx 0.86\,\zeta^k$ (cosine $0.995$--$0.997$ with theory), the two
heads' value circuits are antiparallel per token in every seed (cosine
$-0.994$ to $-1.000$), and --- against the strong form of the gauge
prediction --- the per-head decompositions are as seed-stable as the sums:
training selects a canonical even/odd split rather than a random member of
the loss-equivalent family. The gauge dependence realized in practice is
head relabeling and, more importantly, the \emph{training budget}: the
even-lag component learns $\sim 4\times$ slower, so per-head shares double
between 4k and 16k steps at nearly constant loss. Per-head FRA rankings are
checkpoint-dependent even where they are seed-stable, and a measurement
taken too early reports the head-sum constraint itself as broken.

\subsection{The intervention laws: tracking null and quadratic removal}

In the aggregation regime (next section) the concept signal at the readout
is a sum of path contributions that all estimate the same statistic. A cut
on one edge with concept path share $\rho$ at gain $c$ scales the model's
operational concept estimate by
\begin{equation}
\kappa = 1 - c\rho .
\end{equation}
Two laws follow, with different domains of validity.

\textbf{Tracking-null law (robust).} The projection of the model's output
onto the concept direction declines linearly in $c$ and crosses zero at
$c^\ast = 1/\rho$. This is a projection statement: content the cut removes
or injects \emph{orthogonally} to the concept cannot move it. In a
two-layer counting model where the path decomposition is exact, the
logit-space tracking slope is linear with $R^2 = 1.000$ and the null lands
on $1/\rho$ to machine precision, for two different cuts whose $\rho$
differ by a factor $6.4$ ($c^\ast$ predicted $12.41$ and $1.94$ from the
clean decomposition, measured $12.405$ and $1.938$); after the softmax the
behavioral null shifts by $10$--$13\%$.

\textbf{Quadratic removal law (conditional on concept purity).} Because the
clean model is (near-)Bayes-optimal, cross entropy is stationary at
$\kappa = 1$, so moving the prediction along the clean-to-prior line costs
\begin{equation}\label{quadlaw}
\mathrm{CE}(\kappa) - \mathrm{CE}(1) \approx \tfrac12 C (1-\kappa)^2
\qquad\Longrightarrow\qquad
\RF(c) \approx (c\rho)^2 ,
\end{equation}
where $\RF$ is the removal fraction normalized so that predicting the prior
scores $1$. The hidden premise is that the severed content is a (near-)pure
image of the concept, so that the cut moves the prediction \emph{along}
that line. When it is --- fra\_hmm\_toy's cuts sever concept-selected SAE
latents whose removed signal is $87\%$ concept-aligned --- the law holds
with the stated coefficient. When the cut severs a whole path carrying
mostly other content (in the counting model, paths are only $\approx 30\%$
concept by variance), cross entropy still grows quadratically but with
coefficient $\gamma \gg \rho^2$ ($\approx 0.2$--$0.3$ pointwise versus
$0.0065$), and no gain
returns the model to prior-level loss: the tracking null survives, the CE
null does not.

Three consequences within the purity domain. \textbf{Severing under-reaches
quadratically}: exact severing ($c=1$) of a quarter-share path removes
$6\%$ of the behavior --- a mathematical necessity near optimality rather
than a defect of the attribution. \textbf{The gain-tuned null} sits at
$c^\ast = 1/\rho$. \textbf{One-shot calibration}: a single severing
measurement predicts the null, $c^\ast = 1/\sqrt{\RF(1)}$ --- valid
precisely when the severed content is concept-pure.

On the existing fra\_hmm\_toy gain sweeps a single fitted $\rho = 0.242$
reproduces the whole dose-response ($c^\ast_{\mathrm{pred}} = 4.13$ versus
measured $\approx 4.0$; $\RF(1)$ predicted $0.059$ versus measured $0.062$;
seed 43: $\rho = 0.203$, $c^\ast = 4.93$ versus a measured tracking null
$\approx 4.6$ --- a $7\%$ error, the least accurate prediction in the set), and the
one-shot calibration gives $c^\ast = 1/\sqrt{0.062} = 4.03$. A parameter-free
consistency check follows from the same picture: the removal fraction and
the tracking $R^2$ are two functions of the single scalar $\kappa$, and
their sum equals the clean tracking $R^2$ pointwise --- measured
$0.994$--$1.009$ for $c \le 4$ in the reference sweep (past the null the
identity breaks, $1.17$ at $c=6$), and resolved most cleanly at
concentration 30, where the clean tracking is far from $1$: sums
$0.941$--$0.945$ against clean $0.943$ (seed 43 drifts to $0.958$--$0.994$).
The flatness at the \emph{clean}
value rather than at $1$ carries mechanistic information: the cut contracts
the model's entire operational estimate --- signal and estimation noise
together --- because it is proportional to the model's own latent activations.

Honest deviations: the pointwise $\rho$ implied by individual sweep points
varies by $\pm 10\%$ ($0.232$--$0.277$); past the null the response is
super-quadratic ($\RF(6) = 2.52$ versus $2.10$ predicted); and the null
drifts with the concept prior ($c^\ast \approx 3.4$ to $4.2$ across
Dirichlet concentrations 30 to 0.5), the signature of a locally-linear
cancellation of a mildly nonlinear latent response. Naive ex-ante estimates
of $\rho$ from clean-run projections overpredict $c^\ast$ by up to $38\%$
because the severed signal is only $87\%$ aligned with the concept
direction; the behaviorally correct $\rho$ is recovered by the $c=1$
severing measurement itself.

\section{The mixture regime: flat-profile optimality for exchangeable data}

The fra\_hmm\_toy concept $\omegab$ is a per-sequence mixture weight: given
$\omegab$, tokens are (block-)i.i.d. This makes the token process
\emph{exchangeable}, and exchangeability is exactly the $\lambda \to 1$
degeneration of the spectral machinery: all skip-bigram statistics become
lag-independent,
\begin{equation}
\Btau = \Bb := E_{\omegab}\bigl[\pb(\omegab)\,\pb(\omegab)^\top\bigr]
\quad (\tau \ge 1), \qquad \Db = \mathrm{diag}\bigl(E[\pb(\omegab)]\bigr).
\end{equation}
Within the one-layer MSE framework of the one-layer-transformer note (free
row-stochastic Toeplitz attention profile $\alpha$, effective direct and
context weights optimized out), all attention dependence collapses. Below,
$\lattn$ denotes the attention \emph{benefit} --- the term attention
subtracts from the attention-free loss --- so larger is better:

\begin{theorem}[flat-profile optimality]\label{thm:flat}
For an exchangeable token process, with $\alpha_0 = \alpha(0)$,
$\gamma_0 = \sum_\tau \alpha(\tau)^2$, and $\Rb_0 = \Bb - \Bb\Db^{-1}\Bb
\succeq 0$,
\begin{equation}\label{lattntheta}
\lattn = \tfrac12\, \tr\Bigl( \Rb_0 \bigl[\theta\,(\Db - \Bb) + \Rb_0
\bigr]^{+} \Rb_0 \Bigr),
\qquad
\theta = \frac{\gamma_0 - \alpha_0^2}{(1-\alpha_0)^2},
\end{equation}
a strictly decreasing function of $\theta$ alone. The scalar $\theta$ is the
inverse participation ratio of the off-diagonal lag profile; it is minimized,
uniquely, by the flat profile ($\theta = 1/L$ over $L$ available lags), for
every value of the diagonal share $\alpha_0$, which drops out exactly. As
$L \to \infty$, $\lattn \uparrow \tfrac12 \tr(\Rb_0)$, the total predictive
information the mixture posterior carries about the next token beyond a
single observation.
\end{theorem}

The theorem says three things at once about the aggregation regime.
\textbf{The optimal pattern carries nothing}: no lag gating and no content
gating can improve on the flat profile (the free-$A$ class contains every
QK-realizable pattern), so the concept resides entirely in \emph{what} is
aggregated --- the premise of the QK gauge theorem is derived rather than
assumed. \textbf{The diagonal share is immaterial}, explaining the
$\alpha_0$-dropout observed in the geometric-profile analysis as the
$\lambda \to 1$ limit of a general fact. \textbf{The ceiling
$\tfrac12\tr(\Rb_0)$} prices the concept in MSE units (for the fra\_hmm\_toy
Dirichlet prior, $\tr(\Rb_0) = 0.054$). All identities were verified
independently, including a brute-force $T = 8$ optimization whose argmax is
exactly flat.

For mixture-\emph{with}-grammar data (the actual toy), the skip-bigrams are
a flat mixture floor plus decaying within-component terms, so the optimal
profile superposes a flat component and short-lag geometric structure. The
trained three-layer toy realizes this as head and layer specialization:
both block-0 heads are near-flat (the counting layer, which also carries
$\approx 100\%$ of the concept path share), both block-1 heads are
short-lag geometric, and block~2 has one of each.

The fitted short-lag rates ($0.86$--$0.91$) initially appeared anomalous:
they exceed what was recorded as the dominant grammar eigenvalue ($0.76$),
while the free-$A$ theory predicts optimal rates \emph{below} the process
eigenvalue (on two-state HMMs the closed-form root obeys $\eta < \lambda$
throughout, reproduced by free optimization to $2\times10^{-6}$).
Computing the theory's optimal profile for the \emph{exact} toy process
resolves the anomaly as a double misidentification of the process, not a
failure of the theory: the toy's transition operators carry row mass $1+x$
(a nonstandard normalization, giving true component eigenvalues
$0.85/0.475/0.01$ rather than $1-3x = 0.76$), and mixture components
advance only when selected (lazy chains), so the observable grammar mode
is $\lambda_{\mathrm{eff}} = 1 - \bar\omega_c(1-\zeta_c) = 0.940$. The
MSE-optimal free profile for this process --- computed from exact
skip-bigrams validated against Monte Carlo to $2.6\times10^{-4}$ --- put
through the identical fitting procedure reads $0.8705$, inside the trained
band, and the trained block-1 head profiles lie essentially on the theory
curve. The direction claim survives with the correct eigenvalue:
$0.87 < 0.94$. A methodological byproduct: the optimal profile is a
mixed sum of geometrics ($\{0.978, 0.835\}$ plus a flat floor), and a
single-geometric log-linear fit to such a profile reads high by
construction --- rate-versus-eigenvalue comparisons require both the true
process spectrum and a fit-procedure-matched theory value.

\section{Use versus presence}

Linearize the final-layer residual around the concept coordinates: $\xi =
\mu + \Jb\,\tilde\omegab + \text{noise}$, where $\tilde\omegab = \hat\omegab
- \text{prior}$ has two degrees of freedom, $\Jb = \sum_P \Jb_P$ sums the
per-path concept Jacobians, and the behavioral readout sees $\Wb =
\boldsymbol{U}\Jb$. An FRA-OV cut on edge $e$ at gain $c$ gives $\Jb - c\Kb$
with $\Kb = \sum_{P \ni e} \Jb_P$ confined to the severed subtree's image.

\textbf{Use.} Nulling behavior requires $\boldsymbol{U}(\Jb - c\Kb) = 0$: a
matrix equation with a single unknown $c$, generically unsolvable. It is
solvable in the aggregation regime because every path aggregates the same
statistic, making the path images proportional \emph{as maps} in the readout:
$\boldsymbol{U}\Kb \approx \rho\, \boldsymbol{U}\Jb$. One gain then nulls
all concept components simultaneously --- as observed (all three mixture
components hit the prior together at $c^\ast$). Contrapositive: when a
single gain \emph{can} null a multi-dimensional behavior, that is evidence
the severed paths carry proportional copies --- gain-steering's success is
itself a diagnostic of the aggregation regime.

\textbf{Presence.} Nulling decodability requires $\Jb - c\Kb = 0$ as a full
$d_{\mathrm{model}} \times 2$ matrix: possible only if the severed edge is a
rank bottleneck for the concept. With two or more paths whose residual-space
images are not colinear, no gain achieves it: at the behavioral null the
probe-recoverable presence is untouched (measured $0.985$ against clean
$0.989$, and $0.97$--$0.996$ across every gain, seed, and prior-shift
condition tested). Use and presence come apart exactly because
$\Kb \propto \Jb$ holds in the readout projection but not in the full
residual space.

\textbf{Source cuts and the tax.} A content cut at the residual itself
modifies the common input of all downstream paths and is the only
intervention family that reduces presence --- provided the cut set
\emph{spans} the concept's carrier at that site rather than merely
correlating with it. In the toy the carrier overlaps the token block tags
that the within-component grammar needs, so presence removal is bounded away
from complete by re-aggregation (measured presence floor $0.42$--$0.49$) and
by the entanglement tax: any observer with zero concept information pays at
least $20.8\%$ of the grammar range in collateral (particle-filter bound
from the toy's companion analysis).

\section{Numerical verification}

Every proposition was checked numerically; the table collects the sharpest
check of each. ``Exact'' means agreement at float precision on a direct
construction; predictions marked ex ante were written to the results file
before the measurement ran. All data and scripts live under
\texttt{sprint/code/}, with per-claim JSON outputs. A few quoted quantities
originate outside those outputs and are inherited from the fra\_hmm\_toy
experiment's own artifacts or the working derivations (the $20.8\%$
entanglement tax, the prior-shift null drift $3.4$--$4.2$, the presence
floor $0.42$--$0.49$, the grammar eigenvalue $0.76$, and the two-state
$\eta < \lambda$ scan); they are marked as inherited where used.

\begin{center}
\small
\begin{tabular}{p{4.6cm}p{3.4cm}p{3.4cm}l}
\toprule
claim & predicted & measured & source \\
\midrule
Prop.~\ref{prop:realization}: exact realization & $\ub_d$ exactly &
max error $2.8\times10^{-16}$ & \texttt{phaseA/ideal\_model} \\
Prop.~\ref{prop:warp}: warp in trained model & $w(d) = (1-\zeta^d)^{-1}$ &
RMSE $0.15$ (vs.\ $0.32$ flat), $80\%$ amplitude at $d{=}1$ &
\texttt{phaseA/analyze\_1l} \\
Cor.~\ref{cor:partial}: partial-cut TV at $c{=}1$ & $|c\varphi|m(1{-}m)
= 0.0564$ & $0.0559$; $\Delta$CE to $4$--$5$ decimals &
\texttt{phaseA/fra\_cuts\_1l} \\
Cor.~\ref{cor:asym}: query/key cut asymmetry & query $>$ key &
$2.2\times$ at $c{=}1$ & \texttt{phaseA/fra\_cuts\_1l} \\
gauge-theorem premise, per-sequence & --- & token flip moves weight $56\%$;
gating worth $2.0\times10^{-4}$ nats & \texttt{phaseA3/gating} \\
induction germ at trained optimum & $0$ (S$_3$ symmetry) & tuned gate gains
$6.4\times10^{-7}$ nats & \texttt{phaseA/out/vA3} \\
Thm.~\ref{thm:flat}: flat argmax & flat profile & brute-force $T{=}8$ argmax
exactly flat & \texttt{phaseB/flat\_optimality} \\
quadratic law on existing sweeps & $\RF = (c\rho)^2$, one $\rho$ &
$\rho = 0.242$: $c^\ast$ $4.13$ vs $4.0$, $\RF(1)$ $0.059$ vs $0.062$ &
\texttt{phaseB/fit\_crho} \\
tracking null, ex ante & $c^\ast = 12.41,\ 1.94$ & $12.405,\ 1.938$
(logit space, $R^2 = 1.000$) & \texttt{counting} \\
one-shot calibration & $1/\sqrt{\RF(1)} = 4.03$ & $c^\ast \approx 4.0$ &
\texttt{phaseB/fit\_crho} \\
presence across gains & flat at clean & $0.97$--$0.996$; dip to $0.755$
only at exact severing & \texttt{counting}, sweeps \\
two-head sum constraint ($\zeta = -0.5$) & $\propto \zeta^k$ & cosine
$0.995$--$0.997$, five seeds & \texttt{phaseA2/twohead\_gauge} \\
optimal profile rate, exact process & $0.8705$ (theory, same fit) &
trained $0.86$--$0.91$ & \texttt{phaseC/eta\_discrepancy} \\
SAE repair: content-subspace basis & oracle $\RF = 0.185$ & $0.185$
(FVU $9{\times}10^{-5}$) & \texttt{phaseA4/sae\_repair} \\
\bottomrule
\end{tabular}
\end{center}

\section{Limitations and open gaps}

\begin{enumerate}
\item \textbf{MSE-side theorems, CE-side checks.} The flat-profile theorem
and the free-$A$ machinery are exact in the MSE framework; the
cross-entropy versions of all claims are verified numerically but not
proved. The gauge theorem itself is objective-free (it constrains any
token-independent softmax pattern), but \emph{which} pattern is optimal is
an MSE statement.
\item \textbf{The trained model beats the theory family.} The one-layer
model's CE ($1.07966$) beats the entire gain-scaled constrained-update
family (best $1.08023$). The dissection attributes $64\%$ of the residual
edge to positional-profile freedom and $36\%$ to content-gating, of which
an optimal token-pair gate captures only a quarter; the remaining
position-dependent gating structure is measured but not characterized.
\item \textbf{Quadratic-law scope.} The coefficient $(c\rho)^2$ requires
concept-pure severed content; whole-path cuts obey a quadratic law with a
much larger coefficient and no CE null. Pointwise $\rho$ drifts $\pm 10\%$
within a sweep; the response is super-quadratic past the null; the null
drifts with the concept prior ($3.4$--$4.2$ across concentrations) --- the
gain-tuned null is a calibrated equilibrium, not a structural removal.
\item \textbf{Basis dependence.} All FRA conclusions are conditional on
basis quality. The repair experiments (\S 3) demonstrate a sufficient
training condition (position-demeaned inputs) and the right acceptance
statistic (token-$R^2$ of the summed cut feature) on this toy, but we have
no theorem that these transfer beyond it, and the failure of capacity
scaling to fix entanglement is a single data point.
\item \textbf{Resolved during the sprint: the rate gap.} The apparent
excess of trained pattern rates ($0.86$--$0.91$) over the grammar
eigenvalue ($0.76$) dissolved once the exact process was used (\S 5): the
optimal-profile prediction is $0.87$. What remains is a warning:
the toy's process file normalizes its transition operators nonstandardly,
so quantities derived from the textbook Mess3 formulas (like
$\zeta = 1-3x$) silently misdescribe it.
\item \textbf{Aggregation-regime QK inertness is a weighted statement.}
The raw correlation between QK-cut score deltas and counted tags is not
small ($0.25$); the operative null is the attention-weighted covariance
($\approx 0.02$ against nat-scale deltas). Score-mass screens cannot
distinguish gauge from gating from robust-aggregation nulls; only causal
probes can.
\item \textbf{Per-head attributions are checkpoint-dependent} (\S 4.1):
the even-lag component trains $\sim4\times$ slower, so per-head shares move
at nearly constant loss even when seed-stable.
\end{enumerate}

\end{document}
